\SourcePage{756}{759}
\section{Effective stopping}

In applications of collision theory, the mean energy lost by a colliding
particle is an important quantity to calculate. It is convenient to
characterize this loss by
\begin{equation}
 d\kappa=\sum_n(E_n-E_0)d\sigma_n,
 \tag{149.1}
\end{equation}

\SourcePage{757}{760}
which we shall call the (differential) \emph{effective stopping}; the sum is,
of course, taken over states of both the discrete and continuous spectra, and
$d\kappa$ refers to scattering into the specified solid-angle element
\footnote{If an electron passes through a gas, scattering by different atoms
occurs independently, and $N\,d\kappa$ ($N$ is the number of atoms per unit
volume of the gas) is the energy lost by the electron per unit path length in
collisions that deflect it into the specified solid-angle element.}.

The general formula for the effective stopping of fast electrons is
\begin{equation}
 d\kappa=8\pi\left(\frac{e^2}{\hbar v}\right)^2
 \sum_n(E_n-E_0)\left|\left\langle n\left|\sum_a e^{-i\mathbf q\mathbf r_a}\right|0\right\rangle\right|^2
 \frac{dq}{q^3}
 \tag{149.2}
\end{equation}
($d\sigma_n$ is taken from (148.9)). As in deriving (148.23), we exclude the
region of extremely small angles and again assume that
$(v_0/v)^2\ll\vartheta\ll1$. Then $q$ does not depend on the amount of energy
transferred, and the sum over $n$ can be evaluated in general form.

The sum is evaluated by means of a \emph{sum rule}, derived as follows. The
matrix elements of a quantity $f$ (a function of the coordinates) and of its
time derivative $\dot f$ are related by
\begin{equation}
 (\dot f)_{0n}=-(i/\hbar)(E_n-E_0)f_{0n}.
 \tag{149.3}
\end{equation}
Hence
\[
\begin{aligned}
 \sum_n(E_n-E_0)|f_{0n}|^2
 &=\sum_n(E_n-E_0)f_{0n}(f_{0n})^*
 =\sum_n(E_n-E_0)f_{0n}(f^+)_{n0}\\
 &=i\hbar\sum_n(\dot f)_{0n}(f^+)_{n0}=i\hbar(\dot f f^+)_{00}.
\end{aligned}
\]
The wave functions of the atom's stationary states may be chosen real. The
matrix elements of the coordinate function $f$ then obey $f_{0n}=f_{n0}$,
while the matrix elements in (149.3) correspondingly obey
$(\dot f)_{0n}=-(\dot f)_{n0}$. The sum under consideration may therefore
also be written
\[
 -i\hbar\sum_n(f^+)_{0n}(\dot f)_{n0}=-i\hbar(f^+\dot f)_{00}.
\]
Taking half the sum of the two expressions gives the required sum rule
\begin{equation}
 \sum_n(E_n-E_0)|f_{0n}|^2=\frac{i\hbar}{2}(\dot f f^+-f^+\dot f)_{00}.
 \tag{149.4}
\end{equation}

\SourcePage{758}{761}
Apply it to $f=\sum_a e^{-i\mathbf q\mathbf r_a}$. By (19.2), its time
derivative is represented by the operator
\[
 \hat{\dot f}=-\frac{\hbar}{2m}\sum_a
 [e^{-i\mathbf q\mathbf r_a}(\mathbf q\boldsymbol\nabla_a)
 +(\mathbf q\boldsymbol\nabla_a)e^{-i\mathbf q\mathbf r_a}].
\]
Direct calculation gives
\[
 \hat{\dot f}\hat f^+-\hat f^+\hat{\dot f}=-\frac{i\hbar}{m}q^2Z.
\]
Substitution in (149.4) yields
\begin{equation}
 \sum_n\frac{2m}{\hbar^2q^2}(E_n-E_0)
 \left|\left\langle n\left|\sum_a e^{-i\mathbf q\mathbf r_a}\right|0\right\rangle\right|^2=Z,
 \tag{149.5}
\end{equation}
which performs the summation we need\footnote{Nowhere in deriving this
relation did we use the fact that the state labelled 0 is the atom's normal
state. The relation therefore holds for any initial state.}.

Thus, for the differential effective stopping we find
\begin{equation}
 d\kappa=4\pi\frac{Ze^4}{mv^2}\frac{dq}{q}
 =\frac{2Ze^4}{mv^2}\frac{d\Omega}{\vartheta^2}.
 \tag{149.6}
\end{equation}
Its range of applicability is specified by
$(v_0/v)^2\ll\vartheta\ll1$, that is,
$v_0/v\ll a_0q\ll v/v_0$.

Next define the total effective stopping $\kappa(q_1)$ for all collisions in
which the momentum transfer does not exceed a certain value $q_1$ such that
$v_0/v\ll a_0q_1\ll v/v_0$:
\begin{equation}
 \kappa(q_1)=\sum_n\int_{q_{\min}}^{q_1}(E_n-E_0)d\sigma_n;
 \tag{149.7}
\end{equation}
$q_{\min}$ is given by (148.11). The integral sign cannot be taken outside
the summation sign, since $q_{\min}$ depends on $n$.

Divide the integration region into two parts, from $q_{\min}$ to $q_0$ and
from $q_0$ to $q_1$, where $q_0$ is a value of $q$ such that
$v_0/v\ll q_0a_0\ll1$. Throughout the integration region from $q_{\min}$ to
$q_0$, expression (148.14) may then be used for $d\sigma_n$:
\[
 \kappa(q_0)=8\pi\left(\frac{e}{\hbar v}\right)^2
 \sum_n|\langle n|d_x|0\rangle|^2(E_n-E_0)
 \int_{q_{\min}}^{q_0}\frac{dq}{q},
\]

\SourcePage{759}{762}
whence
\begin{equation}
 \kappa(q_0)=8\pi\left(\frac{e}{\hbar v}\right)^2
 \sum_n|\langle n|d_x|0\rangle|^2(E_n-E_0)
 \ln\frac{q_0\hbar v}{E_n-E_0}.
 \tag{149.8}
\end{equation}
In the region from $q_0$ to $q_1$, on the other hand, the summation over $n$
may be carried out first. It reduces $d\kappa$ to (149.6), whose integration
over $dq$ gives
\begin{equation}
 \kappa(q_1)-\kappa(q_0)=4\pi\frac{Ze^4}{mv^2}\ln\frac{q_1}{q_0}.
 \tag{149.9}
\end{equation}

To transform these expressions, use the sum rule obtained from (149.4) by
setting
\[
 \hat f=\frac{d_x}{e}=\sum_a x_a,\qquad
 \hat{\dot f}=\frac1m\sum_a\hat p_{xa}.
\]
Commuting $\hat f^+$ with $\hat{\dot f}$ gives ($\hat f^+$ here coincides
with $\hat f$)
$\hat{\dot f}\hat f^+-\hat f^+\hat{\dot f}=-(i\hbar/m)Z$, and hence
\footnote{The same comment made concerning (149.5) applies to this relation.}
\begin{equation}
 \sum_nN_{0n}\equiv\sum_n\frac{2m}{(e\hbar)^2}(E_n-E_0)
 |\langle n|d_x|0\rangle|^2=Z.
 \tag{149.10}
\end{equation}

The quantities $N_{0n}$ are called the \emph{oscillator strengths} of the
corresponding transitions. Introduce a mean atomic energy $I$ by the relation
\begin{equation}
 \ln I=\frac{\sum_nN_{0n}\ln(E_n-E_0)}{\sum_nN_{0n}}
 =\frac1Z\sum_nN_{0n}\ln(E_n-E_0).
 \tag{149.11}
\end{equation}
Using (149.10), formula (149.8) may be rewritten as
$\kappa(q_0)=4\pi Ze^4(mv^2)^{-1}\ln(q_0\hbar v/I)$. Adding (149.9), we
finally obtain
\begin{equation}
 \kappa(q_1)=\frac{4\pi Ze^4}{mv^2}\ln\frac{q_1\hbar v}{I}.
 \tag{149.12}
\end{equation}
This formula contains only one constant characteristic of the given atom
\footnote{For hydrogen,
$I=0.55me^4/\hbar^2=14.9\,\text{eV}$. For heavy atoms, good accuracy may be
expected if the constant $I$ is calculated by the Thomas--Fermi method. It is
easy to establish how values of $I$ calculated in this way depend on $Z$. In
the semiclassical case, the energy-level differences correspond to the
normal frequencies of the particle system. The atom's mean normal frequency
is of order $v_0/a_0$; we may therefore conclude that
$I\sim\hbar v_0/a_0$. In the Thomas--Fermi model, the atomic-electron speeds
depend on $Z$ as $Z^{2/3}$, while the atomic dimensions vary as $Z^{-1/3}$.
Thus we find that $I$ must be proportional to $Z$:
$I=\mathrm{const}\cdot Z$. Experimental data give
$\mathrm{const}\sim10\,\text{eV}$.}.

\SourcePage{760}{763}
Expressing $q_1$ in terms of the scattering angle $\vartheta_1$ by
$q_1=mv\vartheta_1/\hbar$, we obtain the effective stopping for scattering
through all angles $\vartheta\leqslant\vartheta_1$:
\begin{equation}
 \kappa(\vartheta_1)=4\pi\frac{Ze^4}{mv^2}\ln\frac{mv^2\vartheta_1}{I}.
 \tag{149.13}
\end{equation}

If $q_1a_0\gg1$ (that is, $\vartheta_1\gg v_0/v$), $\kappa$ may be
expressed as a function of the greatest energy transferred from the incident
electron to the atom. The preceding section showed that when $qa_0\gg1$ the
atom is ionized, with practically the entire momentum $\hbar\mathbf q$ and
the energy transferred to a single atomic electron. Thus $\hbar q$ and
$\varepsilon$ are related as the momentum and energy of an electron:
$\varepsilon=\hbar^2q^2/2m$. Substituting
$q_1^2=2m\varepsilon_1/\hbar^2$ in (149.12), we obtain the effective stopping
in collisions involving energy transfers $\varepsilon\ll\varepsilon_1$:
\begin{equation}
 \kappa(\varepsilon_1)=\frac{2\pi Ze^4}{mv^2}
 \ln\frac{2m\varepsilon_1v^2}{I^2}.
 \tag{149.14}
\end{equation}

We conclude with the following observation. The energy levels of an atom's
discrete spectrum are associated mainly with excitation of one (outer)
electron; excitation of two electrons already usually requires enough energy
to ionize the atom. Consequently, within the sum of oscillator intensities,
transitions to discrete-spectrum states contribute only a fraction of order
unity, whereas ionizing transitions contribute an amount of order $Z$. It
follows that collisions accompanied by ionization play the principal part in
the stopping by heavy atoms.

\ProblemHeading

Determine the total effective stopping of an electron by a hydrogen atom
($I=0.55$ atomic units); at large energy transfers, the faster of the two
colliding electrons is taken to be the primary electron.

\SolutionHeading

When the primary and secondary electrons acquire comparable energies after a
collision, the exchange effect must be included. Therefore, for stopping with
energy transfer from some value $\varepsilon_1$
($I\ll\varepsilon_1\ll v^2$) to the maximum
$\varepsilon_{\max}=E/2=v^2/4$ (because of our definition of the primary
electron), the cross section (148.17) must be used:
\[
\begin{split}
 \kappa(\varepsilon_{\max})-\kappa(\varepsilon_1)
 &=\frac\pi E\int_{\varepsilon_1}^{E/2}\varepsilon
 \left[\frac1{\varepsilon^2}+\frac1{(E-\varepsilon)^2}
 -\frac1{\varepsilon(E-\varepsilon)}\right]d\varepsilon\\
 &=\frac\pi E\left(\ln\frac{E}{8\varepsilon_1}+1\right).
\end{split}
\]

\SourcePage{761}{764}
Adding (149.14), we obtain (in atomic units)\footnote{For collisions of a
positron with a hydrogen atom there is no exchange effect, and the total
stopping follows simply by substituting
$\varepsilon_{\max}=E=v^2/2$ for $\varepsilon_1$ in (149.14):
$\kappa=(4\pi/v^2)\ln(v^2/0.55)$.}
\[
 \kappa=\frac{4\pi}{v^2}\ln\left(\frac{v^2}{2I}\sqrt{e/2}\right)
 =\frac{4\pi}{v^2}\ln\frac{v^2}{0.94}.
\]
